\section{Conclusion}
\label{sec:conclusion}

We present \scb{}, an iterative benchmark designed to measure the quality of code produced by agents when they are forced to make design decisions. We follow explicit construction principles to create black-box problems with only external contracts described. To measure quality, we introduce structural erosion, complexity concentration, and verbosity, the fraction of code that is redundant. No agent is able to solve any of our \transProblems{} problems end-to-end, while the state-of-the-art solve rate is only \overallBestStrictSolveRate{}\% of \overallTotalCheckpoints{} checkpoints. When forced to iterate, agents consistently degrade the codebase's quality. In our maintained-repository calibration panel, agent checkpoints exhibit higher erosion and verbosity, and agent trajectories accumulate both metrics faster than those of sampled Python repository histories. Our analysis of prompting interventions finds that they can improve initial quality issues, but cannot stop the eroding behavior. \scb{} is a first-of-its-kind iterative evaluation that allows proper investigation of how agents act when given decision autonomy.
